(1) 入射束腰处 \(R_c=\infty\)、\(q_{\rm in}=iz_R\)。薄透镜使 \(q\) 经 \(q_{\rm out}=q_{\rm in}/(1-q_{\rm in}/f)\)，等价 \(1/q_{\rm out}=1/q_{\rm in}-1/f\)，自由空间再传播距离 \(d\) 时 \(q(d)=q_{\rm out}+d\)，所有变量有长度量纲。

(2) 整理得到 \(\boxed{q_{\rm out}=-\frac{fz_R^2}{f^2+z_R^2}+i\frac{f^2z_R}{f^2+z_R^2}}\)。令下游新束腰处 \(\operatorname{Re}q(d_*)=0\)，有 \(d_*=fz_R^2/(f^2+z_R^2)\)，新Rayleigh长 \(z_R'=f^2z_R/(f^2+z_R^2)\)，所以 \(w_0'=w_0f/\sqrt{f^2+z_R^2}\)，均为正。

(3) 在自由空间 \(q(d)=(d-d_*)+iz_R'\)，由复束参数定义得 \(w(d)=w_0'\sqrt{1+[(d-d_*)/z_R']^2}\)。若总光功率为 \(P\)，基模中心强度 \(I(0,d)=2P/(\pi w(d)^2)\)，在 \(d=d_*\) 达到 \(2P/(\pi w_0'^2)\)。较短焦距提升理想峰值却也缩短Rayleigh景深。

(4) \(f\gg z_R\) 时透镜非常弱，\(d_*\simeq z_R^2/f\)、\(w_0'\simeq w_0\)；\(f\ll z_R\) 时 \(d_*\simeq f\)、\(w_0'\simeq w_0 f/z_R\)。若近焦光斑趋于衍射极限而透镜孔径截断原Gaussian尾部，则需用完整衍射积分，不能让理想公式无限缩小束腰。
